% Part: incompleteness
% Chapter: theories-computability
% Section: consis-with-q

\documentclass[../../../include/open-logic-section]{subfiles}

\begin{document}

\olfileid{inc}{tcp}{con}

\olsection{$\Th{Q}$ સાથે સુસંગત સિદ્ધાંતો અનિર્ણેય છે}

નીચેનું પ્રમેય કહે છે કે
$\Th{Q}$ પોતે તો અનિર્ણેય
છે જ, પરંતુ $\Th{Q}$ સાથે
વિરોધ ન ધરાવતો કોઈપણ
સિદ્ધાંત અનિર્ણેય છે.

\begin{thm}
$\Th{T}$ એ અંકગણિતની
ભાષામાં એવો કોઈપણ સિદ્ધાંત
માનો કે જે $\Th{Q}$ સાથે
સુસંગત છે (એટલે
$\Th{T} \cup \Th{Q}$
સુસંગત છે). તો
$\Th{T}$ અનિર્ણેય છે.
\end{thm}

\begin{proof}
યાદ રાખો કે $\Th{Q}$ને
સ્વયંસિદ્ધિઓનો સાન્ત ગણ
$!Q_1$, \dots,
$!Q_8$ છે. આપણે તેમને
એક જ સ્વયંસિદ્ધિ
$!E = !Q_1 \land
\dots \land !Q_8$ વડે
બદલી પણ શકીએ.

માનો કે $\Th{T}$ એ
$\Th{Q}$ સાથે સુસંગત
નિર્ણેય સિદ્ધાંત છે.
આવો ગણ લો:
\[
C = \Setabs{!A}{\Th{T} \Proves !E \lif !A}.
\]
$C$ એ $\Th{Q}$ અને
$\Th{\bar Q}$ને અલગ પાડતો
સંગણનીય ગણ છે એવું બતાવીશું,
જે વિરોધાભાસ છે. પ્રથમ,
જો $!A$ એ $\Th{Q}$માં
હોય, તો $\Th{Q}$ની
સ્વયંસિદ્ધિઓમાંથી $!A$
સાબિત થાય છે; નિગમન
પ્રમેયથી પ્રથમ-ક્રમના
તર્કમાં $!E \lif !A$નું
!!a{derivation} મળે છે.
આથી $!A$ એ~$C$માં છે.

બીજી બાજુ, જો $!A$ એ
$\Th{\bar Q}$માં હોય,
તો પ્રથમ-ક્રમના તર્કમાં
$!E \lif \lnot !A$ની
સાબિતી છે. જો $\Th{T}$
$!E \lif !A$ પણ સાબિત
કરે, તો $\Th{T}$ એ
$\lnot !E$ સાબિત કરે,
અને તેથી $\Th{T}
\cup \Th{Q}$ અસુસંગત બને.
પરંતુ $\Th{T} \cup \Th{Q}$
સુસંગત છે એવી આપણી
ધારણા છે. તેથી $\Th{T}$
$!E \lif !A$ સાબિત કરતું
નથી અને $!A$ એ~$C$માં
નથી.

અમે બતાવ્યું કે $!A$ જો
$\Th{Q}$માં હોય તો $C$માં
છે, અને જો $!A$ એ
$\Th{\bar Q}$માં હોય તો
$\Complement{C}$માં છે.
આથી $C$ સંગણનીય વિભાજક
છે, અને આપણને જોઈતો
વિરોધાભાસ મળે છે.
\end{proof}

આ પ્રમેય ઘણું શક્તિશાળી છે.
દાખલા તરીકે તેનાથી આ
અનુવિધાન મળે છે:

\begin{cor}
  અંકગણિતની ભાષામાં
  પ્રથમ-ક્રમનો તર્ક (એટલે
  $\Setabs{!A}{\text{પ્રથમ-ક્રમના તર્કમાં $!A$ સાબિત થાય છે}}$)
  અનિર્ણેય છે.
\end{cor}

\begin{proof}
પ્રથમ-ક્રમનો તર્ક એ
$\emptyset$નાં નિષ્કર્ષોનો
ગણ છે, અને તે
$\Th{Q}$ સાથે સુસંગત છે.
\end{proof}

\end{document}
